% My experience is mainly in in-context learning and prompt optimization at the reasoning and gradient levels. One direction I am interested in exploring is how to control text generation efficiently and optimize prompts in fewer steps. While my work has not been directly focused on AI safety, I think ideas around controllable generation and steering model behavior are closely related to AI alignment and safety.

\documentclass[10pt]{extarticle}

\usepackage[
  left=0.5in,
  right=0.5in,
  top=0.4in,
  bottom=0.6in
]{geometry}
\usepackage{titlesec}
\usepackage{enumitem}


\pagestyle{empty}
\setlength{\parindent}{0pt}
\setlist[itemize]{noitemsep, topsep=2pt, leftmargin=2em}

\usepackage{xcolor}
\definecolor{linkcyan}{RGB}{0,139,139} % darker cyan
\usepackage[colorlinks=true, urlcolor=linkcyan]{hyperref}

\titleformat{\section}{\large\bfseries}{}{0em}{}[\titlerule]

\begin{document}
\vspace{-1em}
\begin{center}
    {\LARGE \textbf{Aunabil Chakma}}\\
    Tucson, AZ \;\textbullet\; (520) 248-0462 \;\textbullet\; aunabilchakma@arizona.edu \;\textbullet\;
    \href{https://github.com/Aunabil4602}{GitHub} \;\textbullet\;
    \href{https://www.linkedin.com/in/aunabil-chakma-56052119b/}{LinkedIn} \;\textbullet\;
    \href{https://aunabil4602.github.io/}{Personal Site} \;\textbullet\;
    \href{https://scholar.google.com/citations?hl=en&user=LQepJpUAAAAJ&view_op=list_works&sortby=pubdate}{Google Scholar}
\end{center}

\vspace{-1em}

\section*{Research Interest}
LLMs, In-Context Learning, Prompt Optimization, Model Alignments

% Interested in NLP, LLMs, ICL, and prompt optimization; experienced in improving model performance via example retrieval and prompt design; also interested in agentic systems and RL-based fine-tuning.

% My research focuses on large language models, especially in-context learning and prompt optimization. My work looks at how better example retrieval, prompt design, and training strategies can improve model performance on NLP tasks such as few-shot relation extraction. I have a paper accepted at ACL (main conference) and another at an EMNLP workshop.

% I have trained and fine-tuned NLP models (e.g., machine translation and sentiment analysis) using supervised learning and prompt-based methods with PyTorch and Hugging Face.

% Yes, my software engineering experience makes me confident in my programming ability and in shipping code to production.

% I am a CS PhD student working on large language models, especially prompt optimization and evaluating model behavior. I am interested in Latitude because this role matches my experience, and I would like to build benchmarks to better understand and improve LLMs.

% My biggest achievement so far is getting a paper accepted at the ACL 2026 Main Conference as my first submission during my PhD, which reflects the impact and quality of my research work.

% Supported CSC445(Algorithms) and CSC450(Algorithms for Bioinformatics) through grading and tutoring.

% Developed prompting strategies for narrative rewriting to improve output quality.
% Conducted the research under Prof. Eduardo Blanco and Prof. Joshua Garland.

% I am a strong fit for Cohere because my PhD research focuses on LLMs, in-context learning, prompt optimization, and NLP, including an accepted ACL 2026 Main Conference paper on few-shot relation extraction. I also bring 3.5+ years of production software engineering experience building scalable Java SaaS systems, combining research depth with real-world deployment experience.

\section*{Education}
\textbf{University of Arizona} \hfill Tucson, AZ\\
\hspace*{1em}Ph.D. in Computer Science \hfill Aug 2024 -- Present

\textbf{Bangladesh University of Engineering and Technology (BUET)} \hfill Dhaka, Bangladesh\\
\hspace*{1em}B.S. in Computer Science and Engineering \hfill Feb 2015 -- Apr 2019

\section*{Selected Publications}
\begin{itemize}
    \item \textbf{Aunabil Chakma}, Mihai Surdeanu, Eduardo Blanco. \textit{Structured Semantic Information Helps Retrieve Better Examples for In-Context Learning Applied to Few-Shot Relation Extraction}. \textbf{ACL 2026 (Oral)}.
    \item \textbf{Aunabil Chakma}, Mihai Surdeanu, Eduardo Blanco. \textit{Two-Stage Prompt Optimization for Few-Shot Relation Extraction: From Reasoning-Guided Search to Gradient-Guided Refinement}. \textbf{ARR May 2026 (Under Review)}.
    % This project improves few-shot relation extraction by turning a single example into a small set of helpful examples for in-context learning. It combines LLM-generated examples with those retrieved via deeper syntactic-semantic similarity, leading to better, more diverse predictions across datasets and language models. This work has been accepted to the ACL 2026 Main Conference.
    \item \textbf{Aunabil Chakma}, Aditya Chakma, Masum Hasan, Soham Khisa, Chumui Tripura, Rifat Shahriyar. \textit{ChakmaNMT: Machine Translation for a Low-Resource and Endangered Language via Transliteration}. \textbf{arXiv}, 2025.
    \item \textbf{Aunabil Chakma}, Masum Hasan. \textit{Leveraging BanglaBERT for Low Resource Sentiment Analysis of Bangla Language}. \textbf{BLP Workshop @ EMNLP2023}.
\end{itemize}

% \section*{Ongoing Research}
% \textbf{Prompt Optimization for LLMs} \hfill Jan 2026 -- Present
% \begin{itemize}
%     \item Exploring textual and gradient-based methods for automatic prompt optimization.
%     \item Working under Prof. Eduardo Blanco and Prof. Mihai Surdeanu.
% \end{itemize}

% \vspace{-1em}
\section*{Work Experience}
\textbf{Research Assistant, University of Arizona} \hfill Jan 2025 -- Dec 2025
\begin{itemize}
    \item Researching narrative rewriting with prompting strategies; advised by Prof. Eduardo Blanco and Prof. Joshua Garland.
\end{itemize}
% \vspace{1em}
\textbf{Teaching Assistant, University of Arizona} \hfill Fall 2024, Spring 2026
\begin{itemize}
    \item Supported {CSC 445 (Algorithms)} and {CSC 450 (Algorithms for Bioinformatics)} through grading and tutoring.
\end{itemize}
% \vspace{1em}
\textbf{Software Engineer (Associate SWE, SWE I, SWE II), Therapservices} \hfill Nov 2020 -- Jul 2024\\
Dhaka, Bangladesh
\begin{itemize}
    \item Built HIPAA-compliant Java SaaS modules and automated backend data/file-transfer workflows for U.S. government and private organizations serving \textbf{500,000+ individuals}.
\end{itemize}

% Built and maintained secure SaaS and backend systems, including data transfer tools that improved workflow efficiency for large-scale healthcare services.
\vspace{-1em}
\section{Academic Service}
\begin{itemize}
  \item Reviewer: EACL 2026 Demo; AACL-IJCNLP 2025 Demo; NeurIPS 2025 Reliable ML Workshop.
  \item Ad-hoc Reviewer: ARR May 2026, ARR May 2025; EMNLP 2025 Demo.
  % \item Secondary Reviewer: ARR March 2026.
\end{itemize}

\vspace{-1em}
\section*{Skills}
\textbf{Programming:} Python, Java, C++, LaTeX\\
\textbf{ML/AI:} PyTorch, TensorFlow, Hugging Face, OpenNMT-py, spaCy\\
% \textbf{Research:} NLP, LLMs, In-Context Learning, Prompt Optimization, Model Training

\vspace{-1em}
\section*{Achievements}
\begin{itemize}
    \item Received \textbf{TA of the Month} in Dec 2024 and Apr 2026.
    \item \textbf{3rd place} (out of 30 teams), Shared Task on Sentiment Analysis of Bangla Social Media Posts, \textbf{BLP Workshop at EMNLP2023}.
    \item \textbf{4th place} (out of 39 teams), BUET Intra CSE-14 Programming Contest 2016.
    \item Solved \textbf{1000+ programming problems}; Codeforces max rating: \textbf{1706}.
\end{itemize}

\end{document}